% Source PDF: source/普通高考/2012/2012江西文.pdf, page 1, question 15.
% pdfimages -list reports no embedded bitmap; the algorithm flowchart is vector/text artwork.
% Grid evidence: tmp/review_flowchart_2012_jiangxi_liberal/grid/page-1-grid100.png and
% tmp/review_flowchart_2012_jiangxi_liberal/grid/q15-block-grid50.png.
% Nodes: 开始; T=0,k=1; sin(k*pi/2)>sin((k-1)*pi/2)?; a=1; a=0;
% T=T+a; k=k+1; k<6?; 输出T; 结束.
% Directed edges: start -> init -> first test; first test(是) -> a=1;
% first test(否) -> a=0; both assignments -> T update -> k update -> second test;
% second test(是) -> left return -> first-test entry; second test(否) -> output -> finish.
% Assignment branches merge only at the T-update entry; the return arrow is a separate
% right-pointing horizontal segment into the first-test entry, followed by one downward arrow.
\begin{tikzpicture}[
  x=1cm,y=1cm,baseline,
  flowarrow/.style={-{Latex[length=2.1mm,width=1.5mm]},line width=0.45pt},
  nodebox/.style={draw,line width=0.45pt,inner sep=2pt,minimum height=.66cm,
    anchor=center,align=center,execute at begin node={\setlength{\parindent}{0pt}}},
  branchlabel/.style={anchor=center,align=center,inner sep=0pt,
    execute at begin node={\setlength{\parindent}{0pt}}},
  term/.style={nodebox,rounded corners=2.5pt,minimum width=1.30cm,text width=1.30cm},
  decision/.style={nodebox,diamond,aspect=2.9,minimum width=6.10cm,minimum height=1.18cm,inner sep=1pt},
  assign/.style={nodebox,minimum width=1.65cm,text width=1.65cm},
  output/.style={nodebox,shape=trapezium,trapezium left angle=70,trapezium right angle=110,
    minimum width=2.10cm,text width=1.85cm,minimum height=.74cm}
]
  % Tight margin around the source flowchart; no node text is indented.
  \path[use as bounding box] (-4.50,-1.75) rectangle (5.95,10.75);

  \node[term] (start) at (0,10.35) {开始};
  \node[nodebox,minimum width=3.20cm,text width=3.20cm] (init) at (0,9.10) {$T=0, k=1$};
  \node[decision] (test1) at (0,7.35)
    {$\sin\dfrac{k\pi}{2}>\sin\dfrac{(k-1)\pi}{2}$?};
  \node[assign] (aone) at (0,5.80) {$a=1$};
  \node[assign] (azero) at (4.55,5.80) {$a=0$};
  \node[nodebox,minimum width=2.75cm,text width=2.75cm] (sum) at (0,4.35) {$T=T+a$};
  \node[nodebox,minimum width=2.35cm,text width=2.35cm] (inc) at (0,2.95) {$k=k+1$};
  \node[decision,minimum width=2.95cm,minimum height=1.02cm] (test2) at (0,1.45) {$k<6$?};
  \node[output] (out) at (0,-0.10) {输出 $T$};
  \node[term] (finish) at (0,-1.35) {结束};

  \coordinate (firstjoin) at (0,8.35);
  \coordinate (assignjoin) at (0,5.05);
  \coordinate (assignright) at (4.55,5.05);
  \coordinate (returnleft) at (-4.00,1.45);
  \coordinate (returntop) at (-4.00,8.35);

  \draw[flowarrow] (start.south) -- (init.north);
  \draw (init.south) -- (firstjoin);
  \draw[flowarrow] (firstjoin) -- (test1.north);

  % First decision branches remain independent until the assignment merge.
  \draw[flowarrow] (test1.south) -- node[right=2pt,pos=.20,yshift=2pt] {是} (aone.north);
  \draw[flowarrow] (test1.east) -- node[above=2pt] {否} (4.55,7.35) -- (azero.north);
  \draw (aone.south) -- (assignjoin);
  \draw[flowarrow] (assignjoin) -- (sum.north);
  \draw[flowarrow] (azero.south) -- (assignright) -- (assignjoin);

  \draw[flowarrow] (sum.south) -- (inc.north);
  \draw[flowarrow] (inc.south) -- (test2.north);

  % Loop return from k<6? (是) enters the first-test junction via one right arrow.
  \draw (test2.west) -- node[above=2pt,pos=.10] {是} (returnleft) -- (returntop);
  \draw[flowarrow] (returntop) -- (firstjoin);

  % Termination branch is independent of the loop.
  \draw[flowarrow] (test2.south) -- node[right=2pt] {否} (out.north);
  \draw[flowarrow] (out.south) -- (finish.north);

\end{tikzpicture}
